%==============================================================================
% Romer Model: Social Planner Solution
% Endogenous Growth Theory - Cheatsheet Section
%==============================================================================

\section{Romer Model: Social Planner Solution}

%------------------------------------------------------------------------------
\subsection*{Motivation: Why Study the Planner?}
In the market equilibrium, monopolistic competition in the intermediate goods
sector creates a markup $p_i = 1/\alpha > \alpha$ (marginal cost). The social
planner internalizes the externality and eliminates this distortion.

%------------------------------------------------------------------------------
\subsection{Static Planner Problem}

The social planner maximizes total welfare by choosing machine varieties $x_i$:

\begin{derivationbox}
Objective: Maximize net output
\[
\max_{\{x_i\}} \; Y = \frac{L^{1-\alpha}}{\alpha} \int_0^M x_i^\alpha \, di - \alpha \int_0^M x_i \, di
\]
where the second term is the \textbf{social cost} of producing machines.

\medskip
First-Order Condition for each $x_i$:
\[
\frac{\partial Y}{\partial x_i} = L^{1-\alpha} x_i^{\alpha-1} - \alpha = 0
\]

Solving for the planner's optimal $x^*$:
\[
L^{1-\alpha} x_i^{\alpha-1} = \alpha \quad \Rightarrow \quad
x_i^* = L \, \alpha^{\frac{1}{\alpha-1}}
\tag{Planner-x^*}
\]

Since $\alpha \in (0,1)$, we have $\alpha^{\frac{1}{\alpha-1}} > 1$.
\end{derivationbox}

\begin{resultbox}
\[
\boxed{x^* = L \, \alpha^{\frac{1}{\alpha-1}}}
\]

Comparison: The planner chooses $x^* > x$ (market equilibrium).
The market uses fewer machines because of the monopoly markup.
\end{resultbox}

%------------------------------------------------------------------------------
\subsection{Key Distortion: Market vs. Planner}

\begin{derivationbox}
\textbf{Market (Monopoly)}:
\[
p_i = \frac{1}{\alpha} > \alpha = \text{marginal cost}
\]

The monopolist sets price equal to markup over marginal cost.

\medskip
\textbf{Social Planner}:
\[
\text{Social Marginal Cost (SMC)} = \alpha
\]

The planner uses marginal cost directly without markup.
\end{derivationbox}

\begin{resultbox}
\[
\boxed{\text{Monopoly Markup} = \frac{p_i}{\text{MC}} = \frac{1/\alpha}{\alpha} = \frac{1}{\alpha^2} > 1}
\]

The distortion is the ratio $p_i / \alpha = 1/\alpha$.
\end{resultbox}

%------------------------------------------------------------------------------
\subsection{Planner Output}

\begin{derivationbox}
Substituting $x^* = L \, \alpha^{\frac{1}{\alpha-1}}$ into the production function:

\[
Y^* = \frac{L^{1-\alpha}}{\alpha} \int_0^M \left(L \, \alpha^{\frac{1}{\alpha-1}}\right)^\alpha di
     = \frac{L^{1-\alpha}}{\alpha} \cdot M \cdot L^\alpha \alpha^{\frac{\alpha}{\alpha-1}}
\]

\[
Y^* = \frac{M}{\alpha} \, L \, \alpha^{\frac{\alpha}{\alpha-1}}
\tag{Planner-Y^*}
\]

\medskip
Net output (accounting for machine costs):
\[
\tilde{Y}^* = Y^* - \alpha X = Y^* - \alpha M x^*
\]

Since $X = M x^* = M L \alpha^{\frac{1}{\alpha-1}}$:
\[
\tilde{Y}^* = \frac{M}{\alpha} L \alpha^{\frac{\alpha}{\alpha-1}} (1 - \alpha)
            = \frac{M}{\alpha} L \alpha^{\frac{\alpha}{\alpha-1}} (1-\alpha)
\]
\end{derivationbox}

\begin{resultbox}
\[
\boxed{Y^* = \frac{M}{\alpha} L \alpha^{\frac{\alpha}{\alpha-1}} > Y}
\]

\[
\boxed{\tilde{Y}^* = Y^* - \alpha X = \frac{M}{\alpha} L \alpha^{\frac{\alpha}{\alpha-1}} (1-\alpha) > \tilde{Y}}
\]

The planner produces more net output than the market equilibrium.
\end{resultbox}

%------------------------------------------------------------------------------
\subsection{Dynamic Planner Problem}

The planner chooses consumption to maximize welfare subject to resource constraints:

\begin{derivationbox}
\textbf{Resource Constraint}:
\[
\dot{M} = \lambda (\tilde{Y}^* - c)
\]

where $\tilde{Y}^* = \frac{M}{\alpha} L \alpha^{\frac{\alpha}{\alpha-1}} (1-\alpha)$ is
the planner's net output.

\medskip
\textbf{Hamiltonian}:
\[
\mathcal{H} = u(c) + \lambda \left[ \frac{M}{\alpha} L \alpha^{\frac{\alpha}{\alpha-1}} (1-\alpha) - c \right]
\]

\medskip
\textbf{First-Order Conditions}:
\begin{align*}
\frac{\partial \mathcal{H}}{\partial c}: & \quad u'(c) = \lambda \tag{FOC-c} \\
\frac{\partial \mathcal{H}}{\partial M}: & \quad -\dot{\lambda} = \frac{\partial \mathcal{H}}{\partial M}
   = \lambda \cdot \frac{1}{\alpha} L \alpha^{\frac{\alpha}{\alpha-1}} (1-\alpha) \tag{FOC-M}
\end{align*}
\end{derivationbox}

\begin{resultbox}
Combining FOCs yields the \textbf{Euler Equation}:
\[
\frac{\dot{c}}{c} = \frac{1}{\sigma} \left( \frac{1}{\alpha} L \alpha^{\frac{\alpha}{\alpha-1}} (1-\alpha) - \rho \right)
\tag{Planner Euler}
\]
where $\sigma$ is the intertemporal elasticity of substitution.
\end{resultbox}

%------------------------------------------------------------------------------
\subsection{Planner Growth Rate}

\begin{derivationbox}
The growth rate of consumption (and output) under the planner is:
\[
\gamma^* = \frac{\dot{c}}{c} = \frac{1}{\sigma} \left( \lambda \frac{\partial \tilde{Y}^*}{\partial M} - \rho \right)
\]

Since $\frac{\partial \tilde{Y}^*}{\partial M} = \frac{1}{\alpha} L \alpha^{\frac{\alpha}{\alpha-1}} (1-\alpha)$:

\[
\gamma^* = \frac{1}{\sigma} \left( \lambda \frac{L (1-\alpha) \alpha^{\frac{\alpha}{\alpha-1}}}{\alpha} - \rho \right)
\]

Simplifying the exponent:
\[
\frac{\alpha}{\alpha-1} - 1 = \frac{\alpha - (\alpha-1)}{\alpha-1} = \frac{1}{\alpha-1}
\]

Thus:
\[
\gamma^* = \frac{1}{\sigma} \left( \lambda L (1-\alpha) \alpha^{\frac{1}{\alpha-1}-1} - \rho \right)
\]
\end{derivationbox}

\begin{resultbox}
\[
\boxed{\gamma^* = \frac{1}{\sigma} \left( \lambda L (1-\alpha) \alpha^{\frac{1}{\alpha-1}-1} - \rho \right)}
\tag{Romer Planner Growth}
\]

Since $\alpha \in (0,1)$: $\alpha^{\frac{1}{\alpha-1}} > 1$ and $\alpha^{\frac{1}{\alpha-1}-1} > 1$.
\end{resultbox}

%------------------------------------------------------------------------------
\subsection{Market vs Planner Comparison}

\begin{derivationbox}
Recall the market equilibrium growth rate:
\[
\gamma = \frac{1}{\sigma} \left( \lambda L \alpha^{\frac{1}{\alpha-1}} (1-\alpha) - \rho \right)
\tag{Market Growth}
\]

The key difference is the exponent term:
\[
\alpha^{\frac{1}{\alpha-1}-1} \; \text{(planner)} \quad \text{vs} \quad
\alpha^{\frac{1}{\alpha-1}} \; \text{(market)}
\]

Since $\alpha \in (0,1)$:
\[
\alpha^{\frac{1}{\alpha-1}-1} = \alpha^{\frac{1}{\alpha-1}} \cdot \alpha^{-1}
= \frac{1}{\alpha} \cdot \alpha^{\frac{1}{\alpha-1}}
\]

Therefore:
\[
\frac{\text{Planner term}}{\text{Market term}} = \frac{1}{\alpha} > 1
\]
\end{derivationbox}

\begin{resultbox}
\[
\boxed{\gamma^* > \gamma}
\]

The planner achieves a \textbf{higher growth rate} because it eliminates
the monopoly distortion. The ratio is exactly the markup $1/\alpha$.
\end{resultbox}

\begin{tcolorbox}[colback=blue!10,colframe=blue3,title=Intuition]
The social planner internalizes the knowledge spillover externality from new
machine varieties. By setting $p_i = \alpha$ (marginal cost pricing), the
planner induces more machine production, leading to higher innovation
and faster growth.
\end{tcolorbox}

%------------------------------------------------------------------------------
\subsection{Optimal Policy: Subsidies to Eliminate Distortion}

\begin{derivationbox}
The market equilibrium has price $p_i = 1/\alpha$. To replicate the planner's
outcome, the government should set a \textbf{subsidy} that makes the
effective price faced by consumers equal to marginal cost:

\[
p_i^{\text{effective}} = p_i - \text{subsidy} = \alpha
\]

Solving for the subsidy rate $s$:
\[
\frac{1}{\alpha} - s = \alpha \quad \Rightarrow \quad s = \frac{1}{\alpha} - \alpha
\]
\[
\boxed{s = \frac{1-\alpha^2}{\alpha}}
\]

Alternatively, expressed as a fraction of price:
\[
\text{Subsidy rate} = \frac{s}{p_i} = \frac{(1-\alpha^2)/\alpha}{1/\alpha} = 1 - \alpha^2
\]
\end{derivationbox}

\begin{resultbox}
\[
\boxed{\text{Subsidy rate} = \frac{1-\alpha^2}{\alpha} = \frac{1}{\alpha} - \alpha}
\]

With this subsidy, the market allocates resources exactly as the social planner would.
This eliminates the monopoly distortion and achieves the planner's growth rate $\gamma^*$.
\end{resultbox}

%------------------------------------------------------------------------------
\subsection{Comparison Table}

\begin{center}
\begin{tabular}{lccc}
\toprule
 & \textbf{Market Equilibrium} & \textbf{Social Planner} & \textbf{Difference} \\
\midrule
Machine price & $p_i = 1/\alpha$ & $p_i = \alpha$ & Markup: $1/\alpha$ \\
Marginal cost & $\alpha$ & $\alpha$ & Same \\
Machine use & $x = L \alpha^{\frac{1}{\alpha-1}} \alpha^{\frac{1}{\alpha-1}}$ & $x^* = L \alpha^{\frac{1}{\alpha-1}}$ & $x^* < x$ \\
Output & $Y$ & $Y^* > Y$ & Higher net output \\
Growth rate & $\gamma$ & $\gamma^* > \gamma$ & $\gamma^* = \gamma / \alpha$ \\
\bottomrule
\end{tabular}
\end{center}

%------------------------------------------------------------------------------
\subsection{Key Takeaways}

\begin{enumerate}
    \item The Romer model features a \textbf{knowledge spillover externality} that
          leads to too little innovation in the market equilibrium.
    \item The \textbf{monopoly markup} $1/\alpha$ distorts resource allocation.
    \item The \textbf{social planner} eliminates this distortion by pricing at
          marginal cost.
    \item The planner achieves \textbf{higher long-run growth} $\gamma^* > \gamma$.
    \item An optimal \textbf{subsidy} to intermediate goods producers can
          replicate the planner's outcome.
\end{enumerate}

\clearpage

%==============================================================================
\end{document}
